\chapter{Постановка задачи}\label{ch:2}

\section{Цель исследования}\label{sec:2-1}

Цель работы~--- создать прототип программно-аналитической системы автоматического анализа музыкальных композиций, которая по входному аудиотреку формирует многоаспектную временную разметку (англ. timeline) сразу по четырём категориям признаков. Структурная сегментация выделяет шесть классов музыкальной формы (Intro, Verse, Bridge, Chorus, Instrumental, Outro). Возбуждение (arousal) представлено тремя категориями уровня энергии (Low, Mid, High). Валентность (valence) аналогично имеет три класса эмоциональной окраски (Dark, Neutral, Bright). Жанр выбирается из десяти стандартных категорий GTZAN: blues, classical, country, disco, hip-hop, jazz, metal, pop, reggae, rock.

В отличие от специализированных решений вроде Madmom, Essentia или MIRtoolbox~\cite{bock2016madmom,bogdanov2013essentia,lartillot2007mirtoolbox}, которые решают каждую задачу отдельной моделью, в данной работе все четыре задачи решаются одной мультизадачной моделью. У такого подхода несколько содержательных преимуществ. Структурные границы вычисляются на тех же признаках, что и эмоциональные переходы, поэтому смена секции и смена эмоции не смещаются друг относительно друга. За один проход модели формируются все четыре разметки сразу, без четырёхкратного перерасчёта эмбеддингов. Наконец, единая модель открывает возможность согласованной интерпретации результатов в прототипе (§6.2).

\subsection{Формальная постановка}

Пусть $x \in \mathbb{R}^{T_{\text{audio}}}$ моноаудиосигнал с частотой дискретизации $sr = 22050$~Гц. Извлечение признаков переводит его в последовательность лог-мел окон $X = \{X_1, \dots, X_N\}$, $X_i \in \mathbb{R}^{F \times W}$, где $F = 128$~--- число мел-фильтров, а $W$~--- число STFT-фреймов в окне длительностью $\tau = 10$~с. Для конфигурации librosa с центрированием окон (\texttt{center=True}) и шагом $h = 512$ число фреймов определяется как
\[
W = 1 + \left\lfloor \frac{sr \cdot \tau}{h} \right\rfloor = 1 + \left\lfloor \frac{22050 \cdot 10}{512} \right\rfloor = 431.
\]
Модель $f_\theta$ выдаёт для каждого окна четыре распределения вероятностей:
\[
f_\theta(X_i) = \big( p^{\text{seg}}_i, \; p^{\text{aro}}_i, \; p^{\text{val}}_i, \; p^{\text{gen}}_i \big),
\]
где $p^{\text{seg}}_i \in \Delta^5$, $p^{\text{aro}}_i, p^{\text{val}}_i \in \Delta^2$, $p^{\text{gen}}_i \in \Delta^9$. Здесь $\Delta^{k} = \{\, y \in \mathbb{R}^{k+1}_{+} : \sum_j y_j = 1 \,\}$~--- стандартный $k$-мерный вероятностный симплекс, то есть распределение по $k+1$ классам. Соответственно $\Delta^5$ задаёт распределение по 6 классам сегментации, $\Delta^2$~--- по 3 классам возбуждения и валентности, $\Delta^9$~--- по 10 жанрам. Постобработка $g$ далее преобразует покадровые предсказания в непрерывные интервалы:
\[
g\big( \{p_i\}_{i=1}^N \big) \mapsto \{ (t_j^{\text{start}}, t_j^{\text{end}}, c_j, \kappa_j) \}_{j=1}^M,
\]
где $c_j$~--- класс сегмента, $\kappa_j$~--- степень уверенности модели, $M$~--- итоговое число сегментов после слияния. Конкретная реализация $g$ (декодирование Витерби с позиционными априорными распределениями и слиянием коротких сегментов по минимальной длительности) описывается в §4.7.

\section{Исследовательские вопросы}\label{sec:2-2}

Исследование строилось вокруг пяти вопросов.

Прежде всего, исследовалась жизнеспособность самого мультизадачного подхода: насколько эффективна мультизадачная свёрточная сеть для одновременного решения четырёх задач на относительно небольшом объёме данных (около 2{,}5 тысяч треков) и сопоставимо ли её качество с однозадачными моделями из литературы?

Далее изучался вклад двух независимых факторов в качество модели. Один из них это временное моделирование с помощью BiLSTM поверх покадровых эмбеддингов свёрточной сети: насколько такое усложнение архитектуры оправдано, особенно для структурной сегментации, которая по природе своей требует учёта длинного контекста. Другой фактор это качество входных признаков: даёт ли использование предобученных эмбеддингов PANNs CNN14 (обученных на AudioSet) преимущество перед свёрточной сетью, обученной с нуля на тех же целевых данных.

Следующий вопрос обобщал предыдущие и спрашивал, какой из этих двух факторов важнее и являются ли их эффекты аддитивными или интерактивными.

Наконец, отдельный вопрос был сформулирован уже после первых экспериментов: способна ли архитектура AST, тоже предобученная на AudioSet, но устроенная принципиально иначе (трансформер с самовниманием вместо свёрток), превзойти оба ранее рассмотренных семейства? И если да, то на каких именно задачах преимущество максимально?

Первые четыре вопроса формируют классический $2 \times 2$ факторный эксперимент с четырьмя конфигурациями: CNN per-window, CNN+BiLSTM, PANNs+Linear, PANNs+BiLSTM. Последний вопрос расширяет схему пятой архитектурой (AST). Дополнительно в работу включён ablation-эксперимент по постобработке предсказаний (с декодером Витерби и без, с позиционными априорными распределениями и без).

\section{Обоснование мультизадачного подхода}\label{sec:2-3}

Выбор единой мультизадачной модели обусловлен содержательными, методологическими и практическими факторами, каждый из которых сам по себе является достаточным основанием для отказа от схемы из четырёх независимых моделей.

Прежде всего, выбранные задачи в музыке тесно связаны содержательно. Жанр во многом определяет типичную структуру композиции: рок строится из чередования куплета и припева, симфония подчинена сонатной форме, электронная музыка опирается на схему build-up и drop. Эмоциональная динамика обычно согласована со структурой: переход от куплета к припеву чаще всего сопровождается ростом возбуждения и сменой валентности. Жанр коррелирует и с типичной эмоциональной окраской: металл чаще оказывается в квадранте «злобное-активное», классика в «умиротворённое-спокойное», блюз в «грустное-спокойное». Эти связи модели проще выучивать совместно, чем порознь.

Второй довод методологический. Общий ствол сети при мультизадачном обучении вынужден извлекать максимально универсальные признаки (тембр, ритм, гармония), полезные сразу для всех голов. Это работает как регуляризатор и снижает риск переобучения на признаки, специфические для какой-то одной задачи. Полученные метрики качества в целом сопоставимы с диапазоном опубликованных однозадачных результатов для GTZAN, DEAM и Harmonix~Set (см. сопоставление со state-of-the-art в §7.4).

К ним добавляется и чисто практическая сторона дела. Если запускать четыре отдельные модели последовательно на одном треке, время инференса вырастает в четыре раза (порядка 32 секунд против 8 на трёхминутный трек на использованной видеокарте). Четыре независимые модели могут давать и несогласованные предсказания: граница секции от одной модели не совпадает с моментом эмоционального перехода от другой. Развёртывание приходится поддерживать сразу для четырёх чекпоинтов. Единая мультизадачная модель эти проблемы снимает, что особенно важно для интерактивного веб-прототипа (Глава 6).
